\uuid{T0XG}
\titre{Vrai/faux}
\auteur{bodin ; hanani ; mzari}
\organisation{exo7}
\licence{CC BY-NC-SA 4.0}
\datecreate{2020-01-20}
\niveau{L1}
\module{Algèbre}
\chapitre{Logique, ensemble, raisonnement}
\sousChapitre{Logique}
\difficulte{2}
\theme{Logique, vrai, faux}
\exo7qcmid{lille-1.logique.008}

\contenu{
\begin{question}
\qclassification{100.01}
\qtags{motcle=vrai/faux}
Soient $P$ et $Q$ deux assertions. Quelles sont les assertions toujours vraies (que $P$ et $Q$ soient vraies ou fausses)  ?
\begin{answers}
    \bad{$P$ et non($P$)}

    \good{non($P$) ou $P$}

    \bad{non($Q$) ou $P$}

    \good{($P$ ou $Q$) ou ($P$ ou non($Q$))}
\end{answers}
\begin{explanations}
On appelle une tautologie une assertion toujours vraie. C'est par exemple le cas de "non($P$) ou $P$", si $P$ est vraie, l'assertion est vraie, si $P$ est fausse, l'assertion est encore vraie !
\end{explanations}
\end{question}
}
